\clearpage\corrigesChapitre{6}{Compter les entités}
\corrigesSouples
\begin{blocLivre}{cor-6-1}\begin{corrige}{1}{Lire les indices}
\begin{questions}
\item $6+12+6=24$ atomes.
\item $100\times12=1200$ atomes d’hydrogène.
\end{questions}
\end{corrige}\end{blocLivre}
\begin{blocLivre}{cor-6-2}\begin{corrige}{2}{Masse d’une molécule}
\begin{questions}
\item $m=1,99\times10^{-26}+2\times2,66\times10^{-26}=\qty{7,31e-26}{kg}$.
\item $\qty{7,31e-23}{g}$ car un kilogramme vaut mille grammes.
\end{questions}
\end{corrige}\end{blocLivre}
\begin{blocLivre}{cor-6-3}\begin{corrige}{3}{Un nombre d’atomes}
\begin{questions}
\item $m=\qty{1,06e-3}{kg}$.
\item $N=1,06\times10^{-3}/(1,06\times10^{-25})=\num{1,00e22}$.
\end{questions}
\end{corrige}\end{blocLivre}
\begin{blocLivre}{cor-6-4}\begin{corrige}{4}{Passer à la mole}
\schema[corrige]{mole}
\begin{questions}
\item $n=3,01\times10^{23}/(6,02\times10^{23})=\qty{0,500}{mol}$.
\item $N=nN_A=\num{1,20e24}$ après arrondi à trois chiffres significatifs.
\end{questions}
\end{corrige}\end{blocLivre}
\begin{blocLivre}{cor-6-5}\begin{corrige}{5}{Le dioxygène}
\begin{questions}
\item $m_e=2\times2,66\times10^{-26}=\qty{5,32e-26}{kg}$.
\item $N=5,32\times10^{-4}/(5,32\times10^{-26})=10^{22}$ ; $n=\qty{0,0166}{mol}$.
\end{questions}
\end{corrige}\end{blocLivre}
\begin{blocLivre}{cor-6-6}\begin{corrige}{6}{Atomes et molécules}
\begin{questions}
\item Il y a 0,200 mol d’atomes de C et 0,400 mol d’atomes de O.
\item Non : le total est 0,600 mol d’atomes, chaque molécule contenant trois atomes.
\end{questions}
\end{corrige}\end{blocLivre}
\begin{blocLivre}{cor-6-7}\begin{corrige}{7}{Revenir à la masse}
\begin{questions}
\item $m=Nm_e=2,00\times10^{22}\times3,00\times10^{-26}=\qty{6,00e-4}{kg}$.
\item $m=\qty{0,600}{g}$.
\end{questions}
\end{corrige}\end{blocLivre}
\begin{blocLivre}{cor-6-8}\begin{corrige}{8}{Une espèce dans un mélange}
\begin{questions}
\item $m=0,200\times5,00=\qty{1,00}{g}$.
\item $N=1,00\times10^{-3}/(5,00\times10^{-26})=\num{2,00e22}$.
\end{questions}
\end{corrige}\end{blocLivre}
\begin{blocLivre}{cor-6-9}\begin{corrige}{9}{Même nombre, masses différentes}
\begin{questions}
\item $m_A=\qty{3,00}{g}$ et $m_B=\qty{6,00}{g}$ : B est deux fois plus massif.
\item Les quantités de matière sont égales puisqu’elles correspondent au même nombre d’entités.
\end{questions}
\end{corrige}\end{blocLivre}
\begin{blocLivre}{cor-6-10}\begin{corrige}{10}{Un solide ionique}
\begin{questions}
\item $N_{\mathrm{Ca}}=6,02\times10^{22}$ et $N_{\mathrm{Cl}}=1,204\times10^{23}$.
\item $n_{\mathrm{Ca}}=\qty{0,100}{mol}$ ; $n_{\mathrm{Cl}}=\qty{0,200}{mol}$. Le mot motif évite de parler à tort de molécules dans le solide.
\end{questions}
\end{corrige}\end{blocLivre}
\begin{blocLivre}{cor-6-11}\begin{corrige}{11}{Une très petite quantité}
\begin{questions}
\item $N=3,00\times10^{-5}/(3,00\times10^{-26})=10^{21}$.
\item $10^{21}/10^9=10^{12}$ : le nombre de molécules vaut mille milliards de fois un milliard.
\end{questions}
\end{corrige}\end{blocLivre}
\begin{blocLivre}{cor-6-12}\begin{corrige}{12}{Détecter un facteur mille}
\demarche{Mettre les deux masses dans la même unité avant de diviser.}
\begin{questions}
\item Il a utilisé deux unités différentes pour les masses.
\item $N=1,80\times10^{-3}/(3,00\times10^{-26})=\num{6,00e22}$ ; $n=6,00/6,02\times10^{-1}=\qty{0,0997}{mol}$.
\end{questions}
\retiens{Une erreur entre grammes et kilogrammes produit un facteur mille sur le comptage.}
\end{corrige}\end{blocLivre}

